% Includable independent-review fragment. Requires amsmath, amssymb,
% amsthm and theorem/proposition/corollary environments in the parent.
\section{Finite tuple domains and quotient equality}
\label{sec:finite-quotient-obstruction}

The one-coordinate majorant theorem has a consequence that is stronger
than failure to express ordinary addition as a single numerical term.
It also obstructs interpretations using finite tuples and a nontrivial
quotient equality, provided the interpreted operation is given by a
tuple of terms. The following statement is purely about functions on
the standard natural numbers; it does not require definability of the
domain or the equivalence relation.

The underlying domination property is the established notion of an
\emph{almost unary function} on a countable set; its closure under
composition belongs to the existing clone-theoretic literature
\cite{pinsker}. The present argument uses that property and, for the
recursor application, the stronger availability of primitive-recursive
majorants. We give the finite-quotient result below with a complete
proof and make no literature-wide priority claim for it.

\begin{theorem}[Finite quotient from coordinate bounds]
\label{thm:finite-quotient}
Let $d\geq1$, let $D\subseteq\mathbb N^d$, let $E$ be an equivalence
relation on $D$, and let $e\in D$. Suppose
$f:D\times D\longrightarrow D$ has a two-sided identity modulo $E$:
\[
                  f(e,z)\mathrel E z,
                  \qquad f(z,e)\mathrel E z
                  \quad(z\in D).
\]
For each output coordinate $i$, suppose there are an input coordinate
$j_i\in\{1,\ldots,d\}$ and a nondecreasing unary function $g_i$ such
that one of the following holds uniformly on $D\times D$:
\[
 f(x,y)_i\leq g_i(x_{j_i}),
 \qquad\hbox{or}\qquad
 f(x,y)_i\leq g_i(y_{j_i}).
\]
A constant bound is also allowed. Then $D/E$ is finite.

More explicitly, put $L(z)=f(e,z)$ and $R(z)=f(z,e)$. A composition
of at most
\[
                   \ell_d=\left\lceil\log_2(d+1)\right\rceil
\]
maps chosen from $L,R$ sends every element of $D$ to an $E$-equivalent
representative in one fixed finite box
$\prod_{i=1}^d\{0,\ldots,C_i\}$. If the coordinate bounds are
primitive recursive, the constants and the word can be obtained
effectively from those bounds, their selected coordinates, and $e$.
\end{theorem}

\begin{proof}
Under $L$, a coordinate bounded through an $x$ input has the constant
bound $g_i(e_{j_i})$; under $R$ it is bounded by $g_i(z_{j_i})$.
For a coordinate bounded through a $y$ input, these roles are reversed.
Thus each coordinate's bound is made independent of the unbounded
input by at least one of the letters $L,R$. Under the other letter it
depends on at most one coordinate. A coordinate with a constant bound
is removed by both letters.

For precision, regard $\{1,\ldots,d\}$ as a finite set of states
tracking dependencies of bounds, not actual functional dependencies.
The transition for a letter sends state $i$ either to the
selected input state $j_i$, or to a terminal symbol $\dagger$ when
that letter gives a constant bound. The terminal symbol is absorbing.
Initially all $d$ states are live. For a current live set $S$, the
number of its states removed by $L$, plus the number removed by $R$,
is at least $|S|$. Choose a letter removing at least
$\lceil|S|/2\rceil$ states, and replace the live set by the images
of the surviving states. Each survivor has only one image, so
\[
                     |S_{\rm next}|\leq\lfloor|S|/2\rfloor.
\]
After at most $\ell_d$ choices the live set is empty, since
$2^{\ell_d}>d$.

The letters are selected from the outside of the composition inward.
If they are $U_1,\ldots,U_k$, the resulting map is
\[
                     W=U_1\circ\cdots\circ U_k;
\]
in an evaluation, $U_k$ is applied first. This order matters for the
dependency calculation. Inductively, each coordinate of the partial
composition is either uniformly bounded, or bounded by a composition
of the unary functions along one live dependency path. Appending the
next letter on the inside replaces that live input coordinate by its
new coordinate bound. Nondecreasingness permits this substitution of
bounds. When the path reaches $\dagger$, the resulting bound is a
constant. The empty final live set therefore gives the constants
$C_1,\ldots,C_d$ asserted in the theorem.

Both $L$ and $R$ map $D$ into $D$. The two identity laws say that
each sends every argument to an $E$-equivalent element. Consequently,
transitivity of $E$ gives $W(z)\mathrel E z$ for every $z\in D$.
No compatibility axiom for $f$ and $E$ is needed for this step.
Every $E$-class thus meets the finite set
\[
               D\cap\prod_{i=1}^d\{0,\ldots,C_i\}.
\]
In particular, the number of classes is at most
$\prod_{i=1}^d(C_i+1)$.
\end{proof}

An arbitrary unary bound may be replaced by its nondecreasing envelope
$g^*(n)=\max_{m\leq n}g(m)$. If $g$ is primitive recursive, so is this
envelope. The monotonicity condition therefore costs no additional
assumption in the application to term majorants. For $d=0$, the
underlying tuple domain has at most one element and the conclusion is
immediate.

\begin{corollary}[No finite-dimensional term interpretation of iPRA]
\label{cor:no-term-interpretation}
Suppose a numerical language has a standard model on $\mathbb N$ in
which every term is bounded by a unary function of one of its free
coordinates, or by a constant. No interpretation of iPRA in this
model can use a fixed finite-dimensional domain, arbitrary quotient
equality, and a term tuple for the interpreted addition operation.
This remains true with fixed numerical parameters in the
interpretation.
\end{corollary}

\begin{proof}
Let $D\subseteq\mathbb N^d$ and $E$ be the proposed interpretation
domain and equality. Choose a representative $e\in D$ of its zero.
The interpreted addition tuple satisfies the hypotheses of
Theorem~\ref{thm:finite-quotient}, because both zero identity laws are
theorems of iPRA. A majorant in a fixed interpretation parameter is a
constant bound and causes no exception. Hence $D/E$ is finite.

On the other hand, iPRA's successor separation and injectivity give
infinitely many distinct standard numerals in every model. Their
interpretations would be infinitely many distinct $E$-classes. This
contradicts the preceding finiteness conclusion.
\end{proof}

Only the two identity laws for addition were used in the finiteness
argument; associativity, commutativity, and cancellation are unnecessary.
Thus the obstruction applies more generally to an interpreted infinite
unital magma. It does not rule out interpretations whose arithmetic
operations are defined by quantified graphs instead of term tuples.
Such relational interpretations are permitted by the original research
question and remain outside this theorem.

\subsection{Sharpness of the composition-length bound}

The logarithmic word-length bound is attained in every positive
dimension. This is a statement about compositions of the two unit
maps, not about optimal proof-checker size or the shortest possible
definition of a representative map.

Fix $d\geq1$, put $k=\lfloor\log_2d\rfloor$, and first suppose
$k\geq1$. Of the $d$ coordinates, label $2^k$ by the words
$w\in\{L,R\}^k$. Write $\sigma(w)$ for the word obtained by removing
its first letter and appending $R$. Take $D=\mathbb N^d$, let $E$
be the universal equivalence relation, and let $e$ be the zero tuple.
For the labeled coordinates define
\[
 f(x,y)_w=
 \begin{cases}
 y_{\sigma(w)},&w\text{ starts with }L,\\
 x_{\sigma(w)},&w\text{ starts with }R.
 \end{cases}
\]
Set every remaining coordinate to zero. The unit laws hold modulo
$E$, and every coordinate is a projection or a constant.

Under a unit-map letter, state $w$ survives exactly when that letter
is its first letter; in that case it moves to $\sigma(w)$. For any
word $\alpha\in\{L,R\}^k$, the state initially labeled $\alpha$
therefore survives all $k$ letters. In the actual composition its
output coordinate is a projection of the input coordinate labeled
$R^k$, so the composition is unbounded. The letters here are again
read from the outside of the composition inward. Hence no word of
length $k$ gives a uniformly bounded map; a shorter bounded word
would remain bounded when extended on the inside and is equally
impossible.

The word $L^{k+1}$ removes every dependency. A state containing $R$
in one of its original $k$ positions is removed on reaching that
position. The only state surviving the first $k$ letters is $L^k$,
which has then become $R^k$ and is removed by the last $L$.
Constant coordinates are removed immediately. Thus the minimum
length is exactly
\[
                 k+1=\left\lceil\log_2(d+1)\right\rceil.
\]

If $k=0$, then $d=1$. Set $f(x,y)=x$, use universal $E$, and take
$e=0$. The map $R$ is the identity and $L$ is constant zero. The
empty word is unbounded and a single $L$ is bounded, giving the
required exact minimum one in this case too.

\subsection{The limits of the obstruction are concrete}

Both identity laws are necessary. On $D=\mathbb N$ with actual
equality, the projection $f(x,y)=x$ has $f(z,e)=z$ for every fixed
$e$ and satisfies the one-coordinate bound, but $D/E$ is infinite.
The omitted identity law is exactly what fails. The opposite
projection illustrates the other one-sided case.

Relational definitions also give a direct escape from the term
hypothesis. In a first-order setting with decidable equality to
zero, the quantifier-free graph
\[
 M(x,y,z)\;\equiv\;
 (x=0\land z=y)\;\lor\;
 (x\ne0\land z=x)
\]
is total and single-valued. It defines the operation that returns
$y$ when $x=0$ and otherwise returns $x$. Its two-sided identity is
zero, so it is an infinite unital magma on $\mathbb N$. Indeed it is
associative: it selects the first nonzero argument in a finite list,
or zero if all are zero. It cannot be represented by an almost unary
numerical term, but its graph is already a finite logical formula.

This example is not an interpretation of iPRA: no addition compiler
or closure under primitive recursion follows from it. It demonstrates
that the remaining relational route is a substantive possibility.
The finite-quotient theorem does not close that route or prove that
the requested smaller axiomatization cannot exist.
